\section{Динамическое программирование \#2}

\subsection{Домашнее задание}
\begin{enumerate}
  \item
    Дан циклический массив, перед $i$-м идёт элемент $(i{-}1) \bmod n$,
    после $i$-го идёт $(i{+}1) \bmod n$. \\
    Рассмотрим операцию \texttt{zero(i) \{ ans += a[next(i)]*a[i]*a[prev(i)]; a[i] = 0 \}}\\
    Найти последовательность операций, максимизирующую \texttt{ans}.
    \begin{enumerate}
      \item $\O(n^3)$.
      \item $\O(n)$.
    \end{enumerate}

  \item
    В Галлии есть $n$ деревень, некоторые из которых соединены дорогами, длины которых --
    натуральные числа. Известно, что от каждой деревни до любой другой можно добраться
    единственным способом. Астерикс и Обеликс решили устроить праздничный забег в честь
    очередной победы над подразделением Цезаря. Они хотят, чтобы забег начинался и заканчивался
    в деревнях, маршрут не проходил ни по какой деревне дважды и имел длину ровно $m$.
    Помогите им посчитать количество различных возможных маршрутов за $\O(nm)$.

  \item
    Пусть даны множества $B, A_1, A_2, \dots , A_m \subset \{0, 1, \dots , n - 1\}$.
    Представьте $B$ в виде объединения минимального количества $A_k$. $\O(m 2^n)$.

  \item
    Найдите во взвешенном графе гамильтонов цикл
    минимального веса, который удовлетворяет дополнительно следующему свойству: сначала
    номера посещенных вершин возрастают, а затем убывают. Время $\O(n^2)$.

\subsection*{Дополнительные задачи}
 \item
    Есть мешок с n камнями с массами $w_i$. Суммарный вес камней равен $W$. Известно, что если столкнуть два камня $i, j$ с весами $w_i \geq w_j$, то они уничтожатся, а на их месте останется камень-осколок весом $w_i - w_j$. Разрешается в произвольном порядке сталкивать по два камня, возвращая результаты столкновений обратно в мешок, до тех пор, пока останется не более одного камня. Определите минимальный возможный вес мешка в конце. $\O(nW)$.

  \item
	Найти максимальное по весу паросочетание на кактусе за
	$\O(n)$. Кактус --- граф, в котором каждое ребро лежит не более чем
	на одном простом цикле.

  \item
	Есть $n$ вещей, у каждой есть стоимость $v_i$ и вес $w_i$. Есть рюкзак, в
	котором можно унести набор вещей суммарного веса не более $W$
	за один подход. За $m = 2^k$ подходов унести вещи максимальной суммарной стоимости.
	Время $\O(3^nk)$.

  \item
	Вам дана доска фанеры размера $n \times m$. В нее было вбито
    несколько гвоздей с целыми координатами (от них остались некрасивые
	дырки). Сколькими способами можно разрезать доску на прямоугольники
	с целыми сторонами так, чтобы ни один из гвоздей не попал внутрь
	прямоугольника. Время: $\O(n^2 4^m)$.

  \item
 	Найти минимальный шаблон такой, что под него не подходит первая строка,
 	но при этом подходит вторая. $\O(n^3)$.

\newpage

\subsection{Решения}

\begin{enumerate}

    \item \dots

    \item
    
    От каждой деревни до другой можно добраться единственным способом, значит, граф является деревом.
    
    Тогда:
    
    \begin{enumerate}
    
    \item Пусть $dp[v][k]$ --- количество путей из вершины $v$ длины $l$.
    \item $dp[v][k] = \sum\limits_{u \in adj[v]} dp[u][k - 1]$, где $adj[v]$ --- соседи вершины $v$.
    \item $dp[v][0] = 1$, так как добраться до вершины длины $0$ можно единственным способом.
    
    \end{enumerate}
    
    Запустим обход в ширину из каждой вершины и будем заполнять массив $dp$. 
    Затем просуммируем все значения $dp[v][m]$ для всех вершин $v$.
    
    \item

    Обозначим $dp[mask]$ как минимальное количество множеств $A_k$, необходимых для покрытия множества,
    заданного битовой маской $mask$.
    
    \begin{enumerate}
    
    \item
    $dp[0] = 0$, так как ничего брать не нужно для остальных масок $dp[mask] = \infty$.

    \item
    Для каждой маски $mask$ из диапазона $[0, 2^n - 1]$ shiftрассмотрим возможность добавить множество $A_k$
    к текущему покрытию.
    
    \item
    Новая маска, соответствующая объединению текущей маски $mask$ и множества $A_k$, будет равна $mask | A_k$.
    
    \item
    Обновляем значение: $dp[new\_mask] = min(dp[new\_mask], dp[mask] + 1)$
    
    \end{enumerate}
    
    Ответ будет находится в $target\_mask$, где $target\_mask$ - битовая маска, соответствующая множеству $B$.
    

    \item \dots

\end{enumerate}

\end{enumerate}
